\subsection{05.05.2026}

\subsubsection{Ziel}

Ziel des Tages war es, die Datenbank mit
echten Messdaten zu testen sowie den Lüfter
anzuschließen und die Steuerung umzusetzen.

\subsubsection{Durchgeführte Arbeiten}

Zunächst wurde überprüft, ob die Messdaten
korrekt in der SQLite-Datenbank gespeichert
werden.

Außerdem wurde die Ansteuerung des Lüfters
vorbereitet und getestet.

Dabei stellte sich heraus, dass der Lüfter
bereits von der vorherigen Gruppe angeschlossen
worden war.

Zusätzlich wurde die Steuerung des Relais
in die bestehende Programmstruktur integriert.

\subsubsection{Problem 1: Fehler in der Datenbank}

Der Zeitstempel wurde nicht korrekt in der
Datenbank gespeichert.

\textbf{Lösung}

Es wurde festgestellt, dass die Funktion
\texttt{datetime("now")} fehlerhaft geschrieben
war.

Nach der Korrektur wurde der Zeitstempel
ordnungsgemäß gespeichert.

\subsubsection{Problem 2: Probleme mit Relais und Lüfter}

Beim Initialisieren der GPIO-Pins konnten
die benötigten Pins nicht korrekt erkannt
werden.

\textbf{Lösung}

Es wurde festgestellt, dass innerhalb des
Projekts unterschiedliche GPIO-Modi verwendet
wurden.

Für die DHT11-Sensoren wurde der Modus
\texttt{GPIO.BCM} genutzt, während beim
Relais versehentlich \texttt{GPIO.BOARD}
verwendet wurde.

Daraufhin wurde das gesamte Projekt auf den
BCM-Modus vereinheitlicht.

Dadurch musste der Lüfteranschluss von
Pin 40 auf GPIO 21 angepasst werden.

Zusätzlich wurde festgestellt, dass zuerst
die DHT11-Sensoren initialisiert werden
müssen, bevor das Relais fehlerfrei gestartet
werden kann.

\subsubsection{Erkenntnisse}

Es wurde deutlich, dass innerhalb eines
GPIO-Projekts ein einheitlicher GPIO-Modus
verwendet werden muss, um Konflikte und
Fehler bei der Pin-Ansteuerung zu vermeiden.

Außerdem wurde bestätigt, dass die
SQLite-Datenbank erfolgreich genutzt werden
kann, um Sensordaten dauerhaft zu speichern.

Zusätzlich wurde erkannt, dass die Reihenfolge
der Hardwareinitialisierung Einfluss auf die
Funktionsfähigkeit einzelner Komponenten haben
kann.

\subsubsection{Nächste Schritte}

Als nächster Schritt soll eine Webseite
erstellt werden, um die Messdaten tabelarisch
darzustellen und die Taupunktlüftung zu
überwachen.